\documentclass[9pt,t]{beamer}
% =====================================================================
% finance-beamer.tex — shared Beamer style for the LLM-in-Finance course
% =====================================================================
% Reusable slide style. Each deck \input's this immediately after
% \documentclass{beamer}. It provides:
%   * the Madrid theme + book colour palette (matches the printed book)
%   * two book-style framing blocks:
%       \begin{bigpicture} ... \end{bigpicture}   ("the bigger picture")
%       \begin{underhood}   ... \end{underhood}    ("under the hood")
%     These mirror the book's `context` and `deepdive` tcolorboxes so the
%     same intuition-vs-mechanism scaffold the reader meets in the book
%     reappears on the slides.
%   * a Python listings style for code frames
%   * appendix conventions: \appendixstart drops a divider and switches
%     to non-numbered backup frames so "extra / less central" material
%     lives after the main talk without inflating the slide count.
%
% Usage:
%   \documentclass{beamer}
%   % =====================================================================
% finance-beamer.tex — shared Beamer style for the LLM-in-Finance course
% =====================================================================
% Reusable slide style. Each deck \input's this immediately after
% \documentclass{beamer}. It provides:
%   * the Madrid theme + book colour palette (matches the printed book)
%   * two book-style framing blocks:
%       \begin{bigpicture} ... \end{bigpicture}   ("the bigger picture")
%       \begin{underhood}   ... \end{underhood}    ("under the hood")
%     These mirror the book's `context` and `deepdive` tcolorboxes so the
%     same intuition-vs-mechanism scaffold the reader meets in the book
%     reappears on the slides.
%   * a Python listings style for code frames
%   * appendix conventions: \appendixstart drops a divider and switches
%     to non-numbered backup frames so "extra / less central" material
%     lives after the main talk without inflating the slide count.
%
% Usage:
%   \documentclass{beamer}
%   % =====================================================================
% finance-beamer.tex — shared Beamer style for the LLM-in-Finance course
% =====================================================================
% Reusable slide style. Each deck \input's this immediately after
% \documentclass{beamer}. It provides:
%   * the Madrid theme + book colour palette (matches the printed book)
%   * two book-style framing blocks:
%       \begin{bigpicture} ... \end{bigpicture}   ("the bigger picture")
%       \begin{underhood}   ... \end{underhood}    ("under the hood")
%     These mirror the book's `context` and `deepdive` tcolorboxes so the
%     same intuition-vs-mechanism scaffold the reader meets in the book
%     reappears on the slides.
%   * a Python listings style for code frames
%   * appendix conventions: \appendixstart drops a divider and switches
%     to non-numbered backup frames so "extra / less central" material
%     lives after the main talk without inflating the slide count.
%
% Usage:
%   \documentclass{beamer}
%   \input{../_style/finance-beamer.tex}
%   \title{...}\subtitle{...}\author{...}\institute{...}\date{\today}
%   \begin{document} ... \end{document}
% =====================================================================

\usetheme{Madrid}
\usecolortheme{whale}
\setbeamertemplate{navigation symbols}{}
\setbeamertemplate{caption}[numbered]
\setbeamercovered{transparent}

% --- Density controls: keep dense, equation-heavy frames on one slide ---
% Slightly smaller body text + tighter lists/blocks. Frame titles and the
% Madrid chrome keep their normal size, so the deck still reads as Madrid,
% just with more vertical room for content.
\setbeamerfont{normal text}{size=\small}
\setbeamertemplate{itemize/enumerate body begin}{\small}
\setbeamertemplate{itemize/enumerate subbody begin}{\footnotesize}
% Tighter block spacing.
\setbeamertemplate{block begin}{%
  \par\vskip2pt%
  \begin{beamercolorbox}[colsep*=.75ex,leftskip=1ex,rightskip=1ex]{block title}%
    \usebeamerfont*{block title}\insertblocktitle%
  \end{beamercolorbox}%
  {\parskip0pt\vskip-1pt}%
  \begin{beamercolorbox}[colsep*=.75ex,vmode,leftskip=1ex,rightskip=1ex]{block body}%
  \vskip1pt}
\setbeamertemplate{block end}{\end{beamercolorbox}\vskip2pt}

\usepackage{amsmath, amssymb}
\usepackage{booktabs}
\usepackage{xcolor}
\usepackage{listings}
\usepackage[skins, breakable]{tcolorbox}

% --- Book colour palette (kept in sync with book/preamble.tex) ---
\definecolor{linkblue}{RGB}{14, 75, 140}
\definecolor{deepdivebg}{RGB}{235, 244, 255}
\definecolor{narrativebg}{RGB}{255, 252, 240}
\definecolor{narrativerule}{RGB}{180, 130, 40}
\definecolor{steelblue}{RGB}{70, 130, 180}
\definecolor{forestgreen}{RGB}{34, 139, 34}
\definecolor{crimson}{RGB}{220, 20, 60}
\definecolor{darkorange}{RGB}{255, 140, 0}
\definecolor{codebg}{RGB}{248, 248, 248}
\definecolor{coderule}{RGB}{210, 210, 210}
\definecolor{keywordblue}{RGB}{0, 100, 180}
\definecolor{commentgray}{RGB}{120, 120, 120}
\definecolor{stringgreen}{RGB}{30, 130, 30}

% Tie the Madrid structural colour to the book's link blue.
\setbeamercolor{structure}{fg=linkblue}
\setbeamercolor{title}{fg=white}
\setbeamercolor{frametitle}{fg=white}

% --- "the bigger picture" : narrative / intuition framing -----------
\newtcolorbox{bigpicture}{%
  enhanced, breakable, boxsep=2pt,
  colback  = narrativebg,
  colframe = narrativerule,
  boxrule  = 0pt, toprule = 1.4pt, bottomrule = 0.4pt,
  leftrule = 0pt, rightrule = 0pt, arc = 0pt,
  left = 6pt, right = 6pt, top = 4pt, bottom = 5pt,
  fonttitle = \footnotesize\scshape\color{narrativerule},
  title = {the bigger picture}, attach title to upper = {\par\smallskip},
}

% --- "under the hood" : the mechanism / technical detail ------------
\newtcolorbox{underhood}{%
  enhanced, breakable, boxsep=2pt,
  colback  = deepdivebg,
  colframe = linkblue,
  boxrule  = 0pt, toprule = 1.4pt, bottomrule = 0.4pt,
  leftrule = 0pt, rightrule = 0pt, arc = 0pt,
  left = 6pt, right = 6pt, top = 4pt, bottom = 5pt,
  fonttitle = \footnotesize\scshape\color{linkblue},
  title = {under the hood}, attach title to upper = {\par\smallskip},
}

% --- Python code style ---------------------------------------------
\lstset{
  basicstyle    = \ttfamily\footnotesize,
  keywordstyle  = \color{keywordblue}\bfseries,
  commentstyle  = \color{commentgray}\itshape,
  stringstyle   = \color{stringgreen},
  showstringspaces = false,
  language      = Python,
  breaklines    = true,
  backgroundcolor = \color{codebg},
  frame         = single,
  rulecolor     = \color{coderule},
  numbers       = none,
  aboveskip     = 4pt, belowskip = 4pt,
}

% --- Appendix conventions ------------------------------------------
% \appendixstart{Title} : begins the "extra material" appendix. Frames
% after it are not counted toward the main deck's frame total, keeping
% the core talk lean while preserving the deeper / optional content.
\newcommand{\appendixstart}[1]{%
  \appendix
  \setbeamertemplate{frametitle continuation}{}%
  \begin{frame}[noframenumbering]
    \begin{center}
      {\Large\scshape\color{linkblue} Appendix}\\[4pt]
      {\large #1}\\[10pt]
      {\footnotesize Extra / optional material — beyond the core lecture.}
    \end{center}
  \end{frame}
}
% Use \begin{frame}[noframenumbering]{...} for each appendix frame.

%   \title{...}\subtitle{...}\author{...}\institute{...}\date{\today}
%   \begin{document} ... \end{document}
% =====================================================================

\usetheme{Madrid}
\usecolortheme{whale}
\setbeamertemplate{navigation symbols}{}
\setbeamertemplate{caption}[numbered]
\setbeamercovered{transparent}

% --- Density controls: keep dense, equation-heavy frames on one slide ---
% Slightly smaller body text + tighter lists/blocks. Frame titles and the
% Madrid chrome keep their normal size, so the deck still reads as Madrid,
% just with more vertical room for content.
\setbeamerfont{normal text}{size=\small}
\setbeamertemplate{itemize/enumerate body begin}{\small}
\setbeamertemplate{itemize/enumerate subbody begin}{\footnotesize}
% Tighter block spacing.
\setbeamertemplate{block begin}{%
  \par\vskip2pt%
  \begin{beamercolorbox}[colsep*=.75ex,leftskip=1ex,rightskip=1ex]{block title}%
    \usebeamerfont*{block title}\insertblocktitle%
  \end{beamercolorbox}%
  {\parskip0pt\vskip-1pt}%
  \begin{beamercolorbox}[colsep*=.75ex,vmode,leftskip=1ex,rightskip=1ex]{block body}%
  \vskip1pt}
\setbeamertemplate{block end}{\end{beamercolorbox}\vskip2pt}

\usepackage{amsmath, amssymb}
\usepackage{booktabs}
\usepackage{xcolor}
\usepackage{listings}
\usepackage[skins, breakable]{tcolorbox}

% --- Book colour palette (kept in sync with book/preamble.tex) ---
\definecolor{linkblue}{RGB}{14, 75, 140}
\definecolor{deepdivebg}{RGB}{235, 244, 255}
\definecolor{narrativebg}{RGB}{255, 252, 240}
\definecolor{narrativerule}{RGB}{180, 130, 40}
\definecolor{steelblue}{RGB}{70, 130, 180}
\definecolor{forestgreen}{RGB}{34, 139, 34}
\definecolor{crimson}{RGB}{220, 20, 60}
\definecolor{darkorange}{RGB}{255, 140, 0}
\definecolor{codebg}{RGB}{248, 248, 248}
\definecolor{coderule}{RGB}{210, 210, 210}
\definecolor{keywordblue}{RGB}{0, 100, 180}
\definecolor{commentgray}{RGB}{120, 120, 120}
\definecolor{stringgreen}{RGB}{30, 130, 30}

% Tie the Madrid structural colour to the book's link blue.
\setbeamercolor{structure}{fg=linkblue}
\setbeamercolor{title}{fg=white}
\setbeamercolor{frametitle}{fg=white}

% --- "the bigger picture" : narrative / intuition framing -----------
\newtcolorbox{bigpicture}{%
  enhanced, breakable, boxsep=2pt,
  colback  = narrativebg,
  colframe = narrativerule,
  boxrule  = 0pt, toprule = 1.4pt, bottomrule = 0.4pt,
  leftrule = 0pt, rightrule = 0pt, arc = 0pt,
  left = 6pt, right = 6pt, top = 4pt, bottom = 5pt,
  fonttitle = \footnotesize\scshape\color{narrativerule},
  title = {the bigger picture}, attach title to upper = {\par\smallskip},
}

% --- "under the hood" : the mechanism / technical detail ------------
\newtcolorbox{underhood}{%
  enhanced, breakable, boxsep=2pt,
  colback  = deepdivebg,
  colframe = linkblue,
  boxrule  = 0pt, toprule = 1.4pt, bottomrule = 0.4pt,
  leftrule = 0pt, rightrule = 0pt, arc = 0pt,
  left = 6pt, right = 6pt, top = 4pt, bottom = 5pt,
  fonttitle = \footnotesize\scshape\color{linkblue},
  title = {under the hood}, attach title to upper = {\par\smallskip},
}

% --- Python code style ---------------------------------------------
\lstset{
  basicstyle    = \ttfamily\footnotesize,
  keywordstyle  = \color{keywordblue}\bfseries,
  commentstyle  = \color{commentgray}\itshape,
  stringstyle   = \color{stringgreen},
  showstringspaces = false,
  language      = Python,
  breaklines    = true,
  backgroundcolor = \color{codebg},
  frame         = single,
  rulecolor     = \color{coderule},
  numbers       = none,
  aboveskip     = 4pt, belowskip = 4pt,
}

% --- Appendix conventions ------------------------------------------
% \appendixstart{Title} : begins the "extra material" appendix. Frames
% after it are not counted toward the main deck's frame total, keeping
% the core talk lean while preserving the deeper / optional content.
\newcommand{\appendixstart}[1]{%
  \appendix
  \setbeamertemplate{frametitle continuation}{}%
  \begin{frame}[noframenumbering]
    \begin{center}
      {\Large\scshape\color{linkblue} Appendix}\\[4pt]
      {\large #1}\\[10pt]
      {\footnotesize Extra / optional material — beyond the core lecture.}
    \end{center}
  \end{frame}
}
% Use \begin{frame}[noframenumbering]{...} for each appendix frame.

%   \title{...}\subtitle{...}\author{...}\institute{...}\date{\today}
%   \begin{document} ... \end{document}
% =====================================================================

\usetheme{Madrid}
\usecolortheme{whale}
\setbeamertemplate{navigation symbols}{}
\setbeamertemplate{caption}[numbered]
\setbeamercovered{transparent}

% --- Density controls: keep dense, equation-heavy frames on one slide ---
% Slightly smaller body text + tighter lists/blocks. Frame titles and the
% Madrid chrome keep their normal size, so the deck still reads as Madrid,
% just with more vertical room for content.
\setbeamerfont{normal text}{size=\small}
\setbeamertemplate{itemize/enumerate body begin}{\small}
\setbeamertemplate{itemize/enumerate subbody begin}{\footnotesize}
% Tighter block spacing.
\setbeamertemplate{block begin}{%
  \par\vskip2pt%
  \begin{beamercolorbox}[colsep*=.75ex,leftskip=1ex,rightskip=1ex]{block title}%
    \usebeamerfont*{block title}\insertblocktitle%
  \end{beamercolorbox}%
  {\parskip0pt\vskip-1pt}%
  \begin{beamercolorbox}[colsep*=.75ex,vmode,leftskip=1ex,rightskip=1ex]{block body}%
  \vskip1pt}
\setbeamertemplate{block end}{\end{beamercolorbox}\vskip2pt}

\usepackage{amsmath, amssymb}
\usepackage{booktabs}
\usepackage{xcolor}
\usepackage{listings}
\usepackage[skins, breakable]{tcolorbox}

% --- Book colour palette (kept in sync with book/preamble.tex) ---
\definecolor{linkblue}{RGB}{14, 75, 140}
\definecolor{deepdivebg}{RGB}{235, 244, 255}
\definecolor{narrativebg}{RGB}{255, 252, 240}
\definecolor{narrativerule}{RGB}{180, 130, 40}
\definecolor{steelblue}{RGB}{70, 130, 180}
\definecolor{forestgreen}{RGB}{34, 139, 34}
\definecolor{crimson}{RGB}{220, 20, 60}
\definecolor{darkorange}{RGB}{255, 140, 0}
\definecolor{codebg}{RGB}{248, 248, 248}
\definecolor{coderule}{RGB}{210, 210, 210}
\definecolor{keywordblue}{RGB}{0, 100, 180}
\definecolor{commentgray}{RGB}{120, 120, 120}
\definecolor{stringgreen}{RGB}{30, 130, 30}

% Tie the Madrid structural colour to the book's link blue.
\setbeamercolor{structure}{fg=linkblue}
\setbeamercolor{title}{fg=white}
\setbeamercolor{frametitle}{fg=white}

% --- "the bigger picture" : narrative / intuition framing -----------
\newtcolorbox{bigpicture}{%
  enhanced, breakable, boxsep=2pt,
  colback  = narrativebg,
  colframe = narrativerule,
  boxrule  = 0pt, toprule = 1.4pt, bottomrule = 0.4pt,
  leftrule = 0pt, rightrule = 0pt, arc = 0pt,
  left = 6pt, right = 6pt, top = 4pt, bottom = 5pt,
  fonttitle = \footnotesize\scshape\color{narrativerule},
  title = {the bigger picture}, attach title to upper = {\par\smallskip},
}

% --- "under the hood" : the mechanism / technical detail ------------
\newtcolorbox{underhood}{%
  enhanced, breakable, boxsep=2pt,
  colback  = deepdivebg,
  colframe = linkblue,
  boxrule  = 0pt, toprule = 1.4pt, bottomrule = 0.4pt,
  leftrule = 0pt, rightrule = 0pt, arc = 0pt,
  left = 6pt, right = 6pt, top = 4pt, bottom = 5pt,
  fonttitle = \footnotesize\scshape\color{linkblue},
  title = {under the hood}, attach title to upper = {\par\smallskip},
}

% --- Python code style ---------------------------------------------
\lstset{
  basicstyle    = \ttfamily\footnotesize,
  keywordstyle  = \color{keywordblue}\bfseries,
  commentstyle  = \color{commentgray}\itshape,
  stringstyle   = \color{stringgreen},
  showstringspaces = false,
  language      = Python,
  breaklines    = true,
  backgroundcolor = \color{codebg},
  frame         = single,
  rulecolor     = \color{coderule},
  numbers       = none,
  aboveskip     = 4pt, belowskip = 4pt,
}

% --- Appendix conventions ------------------------------------------
% \appendixstart{Title} : begins the "extra material" appendix. Frames
% after it are not counted toward the main deck's frame total, keeping
% the core talk lean while preserving the deeper / optional content.
\newcommand{\appendixstart}[1]{%
  \appendix
  \setbeamertemplate{frametitle continuation}{}%
  \begin{frame}[noframenumbering]
    \begin{center}
      {\Large\scshape\color{linkblue} Appendix}\\[4pt]
      {\large #1}\\[10pt]
      {\footnotesize Extra / optional material — beyond the core lecture.}
    \end{center}
  \end{frame}
}
% Use \begin{frame}[noframenumbering]{...} for each appendix frame.


\title{Practical Session 10: Portfolios, Views, and Backtests}
\subtitle{Hands-On --- Efficient Frontier, BL Views, and a Cost-Aware Walk-Forward}
\author{Juan F.\ Imbet}
\institute{Paris Dauphine -- PSL University}
\date{\today}

\begin{document}

\begin{frame}\titlepage\end{frame}

% ============================================================
\begin{frame}{Session Overview}
\textbf{Duration:} 1 hour (50 min problems $+$ 10 min debrief)

\medskip
\textbf{Agenda:}
\begin{enumerate}
  \item \textbf{Recap} (5 min): mean-variance, BL posterior, OOS Sharpe, turnover cost.
  \item \textbf{Problem 1} (12 min): how much to trust an LLM's return estimates --- min-variance vs.\ tangency, and signal-error sensitivity.
  \item \textbf{Problem 2} (12 min): map an LLM sentiment score to a Black--Litterman view.
  \item \textbf{Problem 3} (10 min): does a daily signal survive transaction costs?
  \item \textbf{Case study} (15 min): walk-forward backtest of a sentiment factor in Python
        (\texttt{code/notebooks/10-portfolio-quant-trading/}).
  \item \textbf{Discussion} (6 min): when is an LLM signal worth deploying?
\end{enumerate}

\medskip
\begin{block}{Goal}
Leave able to translate a textual signal into a portfolio weight \emph{and} to judge,
quantitatively, whether the resulting strategy is deployable.
\end{block}
\end{frame}

% ============================================================
\begin{frame}{Recap: The Four Formulas You Will Use}
\begin{underhood}
\textbf{Mean-variance:} $\mu_p=\boldsymbol{\mu}^\top\mathbf{w}$,\;
$\sigma_p^2=\mathbf{w}^\top\boldsymbol{\Sigma}\mathbf{w}$;\quad
tangency $\mathbf{w}^{\text{SR}}\propto\boldsymbol{\Sigma}^{-1}\boldsymbol{\mu}$.
\end{underhood}

\medskip
\textbf{Black--Litterman posterior mean:}
\[
\hat{\boldsymbol{\mu}}=
\bigl[(\tau\boldsymbol{\Sigma})^{-1}+\mathbf{P}^\top\boldsymbol{\Omega}^{-1}\mathbf{P}\bigr]^{-1}
\bigl[(\tau\boldsymbol{\Sigma})^{-1}\boldsymbol{\Pi}+\mathbf{P}^\top\boldsymbol{\Omega}^{-1}\mathbf{q}\bigr],
\quad q_i=\alpha s_i+\Pi_i.
\]

\textbf{Out-of-sample Sharpe (annualised):}
\[
\text{SR}^{\text{OOS}}=\frac{\bar r_p^{\text{OOS}}}{\hat\sigma_p^{\text{OOS}}}\sqrt{252}.
\]

\textbf{Turnover transaction cost:}
\[
\text{TC}_t=c\sum_i|\Delta w_{i,t}|,\qquad \Delta w_{i,t}=w_{i,t}-w_{i,t-1}.
\]
\end{frame}

% ============================================================
\begin{frame}{Problem 1: How Much Should You Trust the LLM's Returns?}
An LLM sentiment model reads each firm's news flow and outputs \textbf{estimated} monthly
excess returns $\boldsymbol{\mu}$ for two assets; the covariance $\boldsymbol{\Sigma}$ is
from price history:
\[
\boldsymbol{\mu}_{\text{LLM}}=\begin{pmatrix}0.8\%\\0.4\%\end{pmatrix},\qquad
\boldsymbol{\Sigma}=\begin{pmatrix}0.0036 & 0.0009\\ 0.0009 & 0.0016\end{pmatrix}.
\]
($\sigma_1=6\%$, $\sigma_2=4\%$, $\rho=0.375$.)

\medskip
\textbf{Tasks:}
\begin{itemize}
  \item[\textbf{A.}] The minimum-variance portfolio \emph{ignores} $\boldsymbol{\mu}$ ---
        the ``trust the LLM nothing'' baseline: $w^\star=\dfrac{\sigma_2^2-\sigma_{12}}{\sigma_1^2+\sigma_2^2-2\sigma_{12}}$.
  \item[\textbf{B.}] The tangency portfolio \emph{fully trusts} the LLM:
        $\mathbf{w}^{\text{SR}}\propto\boldsymbol{\Sigma}^{-1}\boldsymbol{\mu}_{\text{LLM}}$, normalised.
  \item[\textbf{C.}] \textbf{Signal reliability.} Suppose the LLM over-states asset~1 by 50\%
        ($\mu_1:0.8\%\to0.4\%$). Recompute the tangency weights. How far do they move? What does
        this say about sizing positions on a noisy LLM signal vs.\ shrinking toward $1/N$
        (DeMiguel et al., 2009)?
\end{itemize}
\end{frame}

% ============================================================
\begin{frame}{Problem 2: From LLM Score to a BL View}
\textbf{Setup.} An LLM scores firm $i$'s earnings call: $s_i=+0.6$ (bullish).
Equilibrium prior $\Pi_i=0.30\%$ monthly. Calibration constant $\alpha=0.8\%$;
view-noise scale $\sigma_q^2=(0.5\%)^2$.

\medskip
\textbf{Tasks:}
\begin{itemize}
  \item[\textbf{A.}] Compute the view value $q_i=\alpha s_i+\Pi_i$.
  \item[\textbf{B.}] Compute the view uncertainty $\Omega_{ii}=\sigma_q^2(1-|s_i|)$.
  \item[\textbf{C.}] A second firm scores $s_j=+0.05$ (near-neutral).
        Compute $q_j,\ \Omega_{jj}$. Which view will move the posterior more, and why?
  \item[\textbf{D.}] \emph{Conceptual:} you accidentally calibrate $\alpha$ on the same
        window you evaluate on. Name the bias and the fix.
\end{itemize}

\medskip
\begin{block}{Hint}
Posterior pull toward a view scales with view \emph{precision} $\Omega_{ii}^{-1}$.
Strong signal $\Rightarrow$ small $\Omega_{ii}\Rightarrow$ large precision.
\end{block}
\end{frame}

% ============================================================
\begin{frame}{Problem 3: Does the Signal Survive Costs?}
\textbf{Setup.} A daily long-short sentiment strategy with:
\begin{itemize}
  \item Gross (pre-cost) annualised return $4.0\%$, volatility $8.0\%$.
  \item Average \emph{one-way} daily turnover $20\%$; cost $c=0.04\%$ per unit turnover.
  \item 252 trading days.
\end{itemize}

\medskip
\textbf{Tasks:}
\begin{itemize}
  \item[\textbf{A.}] Annualised transaction cost
        $\approx \text{turnover}\times 252\times c$.
  \item[\textbf{B.}] Net annualised return $=$ gross $-$ cost. Net Sharpe (assume vol unchanged).
  \item[\textbf{C.}] By how much must one-way turnover fall to keep net return $\ge 2\%$,
        holding $c$ fixed?
  \item[\textbf{D.}] \emph{Capacity:} the strategy wants \$10M in a stock with \$50M ADV.
        What participation rate is that, and does it breach the $5$--$10\%$ guideline?
\end{itemize}
\end{frame}

% ============================================================
\begin{frame}{Case Study: Walk-Forward Sentiment Backtest}
\textbf{Goal.} Run a leakage-free walk-forward backtest of a language factor and
read off the out-of-sample, cost-adjusted Sharpe.

\medskip
\textbf{Notebook:} \texttt{code/notebooks/10-portfolio-quant-trading/demo.ipynb}\\
\textbf{Frontier figure:} \texttt{gen\_efficient\_frontier.py} (deterministic, no network).

\medskip
\textbf{What you will do:}
\begin{enumerate}
  \item Build quintile language-factor returns from per-firm sentiment scores.
  \item Slide the walk-forward window; \emph{fit calibration on train, evaluate on test}.
  \item Subtract turnover costs; compute $\text{SR}^{\text{OOS}}$.
  \item Answer three diagnostic questions about leakage and decay.
\end{enumerate}

\medskip
\begin{alertblock}{Discipline}
Any parameter ($\alpha$, $\lambda$, $W_{\max}$) touched by the test window invalidates the result.
\end{alertblock}
\end{frame}

% ============================================================
\begin{frame}[fragile]{Case Study: The Code}
\begin{lstlisting}[basicstyle=\ttfamily\scriptsize]
import numpy as np, pandas as pd
def language_factor(scores, fwd_ret):     # firms x dates, scores in [-1,1]
    """Top-minus-bottom quintile return from demeaned sentiment."""
    s = scores.sub(scores.mean(axis=0), axis=1)   # cross-sectional demean
    lf = {}
    for t in s.columns:
        q = s[t].rank(pct=True)
        top, bot = q >= 0.8, q <= 0.2
        lf[t] = fwd_ret.loc[top, t].mean() - fwd_ret.loc[bot, t].mean()
    return pd.Series(lf)

def walk_forward_sharpe(lf, train=252, step=21, turnover=0.15, c=5e-4):
    oos = []
    for m in range(0, len(lf) - train - step, step):
        test = lf.iloc[m + train : m + train + step]  # never seen in fit
        oos.append(test - turnover * c)               # cost-adjust
    r = pd.concat(oos)
    return float(r.mean() / r.std() * np.sqrt(252))
# sharpe = walk_forward_sharpe(language_factor(scores, fwd_ret))
\end{lstlisting}
\end{frame}

% ============================================================
\begin{frame}{Case Study: Diagnostic Questions}
Run the notebook, then answer:

\medskip
\textbf{Q1 (leakage).} The demean step uses \texttt{scores.mean(axis=0)} --- a
\emph{cross-sectional} (per-date) mean. Why is that safe, whereas a \emph{time-series}
mean over all dates would leak?

\medskip
\textbf{Q2 (decay).} Replace raw scores with EWMA-decayed scores
($\lambda\in\{0.1,0.3,0.5\}$). How does $\text{SR}^{\text{OOS}}$ move? Relate to the
1--7 day half-life rule.

\medskip
\textbf{Q3 (costs).} Sweep $c\in\{0,\,5,\,10\}$ bps. At what cost does the
Sharpe cross below the $1.5$ paper-trading threshold? Connect to the capacity constraint.

\medskip
\begin{block}{Optional extension}
Feed the per-firm view $q_i=\alpha s_i+\Pi_i$ into a Black--Litterman optimizer instead of
the raw quintile sort. Does BL's prior shrinkage raise or lower turnover? Why?
\end{block}
\end{frame}

% ============================================================
\begin{frame}{Discussion: Is This Signal Worth Deploying?}
For each scenario, decide \emph{go / no-go} and justify with the lecture's criteria.

\medskip
\begin{enumerate}
  \item A large-cap US equity sentiment signal: OOS Sharpe $1.7$ in-backtest, but
        decays within 2 days and the desk runs \$2B. Crowding? Capacity?
  \item A small-cap emerging-market filing signal: OOS Sharpe $1.2$, low coverage,
        wide spreads (tens of bps). Grossman--Stiglitz niche?
  \item A real-time 8-K risk-alert system: no return claim, but cuts tail losses.
        How would you value it, and audit it under SR 11-7?
\end{enumerate}

\medskip
\begin{block}{Frame your answer around}
OOS Sharpe \emph{and} $T^*$, transaction costs, capacity / ADV, crowding,
auditability, and whether the LLM is acting as an extractor or (wrongly) as an optimizer.
\end{block}
\end{frame}

% ============================================================
\begin{frame}{Solution: Problem 1 --- Trusting the LLM's Returns}
\textbf{A. Minimum-variance} (ignores $\boldsymbol{\mu}$):
$w^\star=\dfrac{0.0016-0.0009}{0.0034}\approx 0.21\Rightarrow \mathbf{w}_{\min}\approx(0.21,\,0.79).$

\medskip
\textbf{B. Tangency} (fully trusts $\boldsymbol{\mu}_{\text{LLM}}$),
$\boldsymbol{\Sigma}^{-1}\boldsymbol{\mu}\propto(9.2,\,6.7)\!\times\!10^{-6}
\Rightarrow \mathbf{w}^{\text{SR}}\approx(0.58,\,0.42).$

\medskip
\textbf{C. Signal reliability.} With $\mu_1$ corrected to $0.4\%$,
$\boldsymbol{\Sigma}^{-1}\boldsymbol{\mu}\propto(2.8,\,10.8)\!\times\!10^{-6}
\Rightarrow \mathbf{w}^{\text{SR}}\approx(0.21,\,0.79)$ --- it collapses all the way to the
min-variance baseline. A \emph{single} 50\% error in the LLM's expected return swings the
allocation by $\sim$37 points. The takeaway: position sizing is brutally sensitive to LLM
signal noise, which is exactly why you shrink toward $1/N$ / use Black--Litterman rather than
plug raw LLM $\boldsymbol{\mu}$ into mean--variance (DeMiguel et al., 2009).
\end{frame}

% ============================================================
\begin{frame}{Solution: Problem 2 --- LLM Score to BL View}
\textbf{Firm $i$ ($s_i=0.6$):}
\[
q_i=0.8\%\times0.6+0.30\%=0.48\%+0.30\%=\mathbf{0.78\%},
\quad
\Omega_{ii}=(0.5\%)^2(1-0.6)=0.4\times2.5\times10^{-5}=\mathbf{1.0\times10^{-5}}.
\]

\textbf{Firm $j$ ($s_j=0.05$):}
\[
q_j=0.8\%\times0.05+0.30\%=\mathbf{0.34\%},
\quad
\Omega_{jj}=(0.5\%)^2(1-0.05)=\mathbf{2.375\times10^{-5}}.
\]

\textbf{C.} View $i$ moves the posterior more: it has both a larger deviation from the
prior ($0.78\%$ vs.\ $0.34\%$) \emph{and} higher precision ($\Omega_{ii}^{-1}$ is
$\approx2.4\times$ larger). The near-neutral signal $j$ lets the prior dominate --- exactly
the intended behaviour of $\Omega_{ii}=\sigma_q^2(1-|s_i|)$.

\medskip
\textbf{D.} Calibrating $\alpha$ on the evaluation window is \emph{model-selection
look-ahead leakage}. Fix: estimate $\alpha$ only on data predating each test window
(walk-forward).
\end{frame}

% ============================================================
\begin{frame}{Solution: Problem 3 --- Survive Costs?}
\textbf{A. Annualised cost:}
\[
0.20\times252\times0.04\%=0.20\times252\times0.0004=\mathbf{2.016\%}\ \text{p.a.}
\]

\textbf{B. Net return / Sharpe:}
\[
\text{net}=4.0\%-2.016\%\approx\mathbf{1.98\%},\quad
\text{net SR}=\frac{1.98\%}{8.0\%}\approx\mathbf{0.25}\ \text{(vs.\ gross }0.50).
\]
Costs \emph{halve} the Sharpe --- a marginal signal is nearly destroyed.

\medskip
\textbf{C. Turnover for net $\ge2\%$:} need cost $\le 2.0\%$, i.e.
\[
\text{turnover}\times252\times0.0004\le0.02
\Rightarrow \text{turnover}\le\frac{0.02}{0.1008}\approx\mathbf{19.8\%}.
\]
Barely below the current $20\%$ --- essentially no slack; the strategy is cost-bound.

\medskip
\textbf{D. Capacity:} \$10M / \$50M ADV $=\mathbf{20\%}$ participation ---
\emph{double} the $5$--$10\%$ guideline. Expect material price impact; scale the book down.
\end{frame}

% ============================================================
\begin{frame}{Wrap-Up: Key Takeaways}
\begin{itemize}
  \item \textbf{P1:} On two assets, $1/N$ already sits near tangency --- the optimization
        prize is small once $\boldsymbol{\mu}$ is estimated.
  \item \textbf{P2:} The map $q_i=\alpha s_i+\Pi_i$, $\Omega_{ii}=\sigma_q^2(1-|s_i|)$
        turns a raw LLM score into a confidence-weighted BL view; neutral signals defer to the prior.
  \item \textbf{P3:} Transaction costs and turnover are first-order --- a 4\% gross
        strategy can net 2\% and lose half its Sharpe; watch ADV participation.
  \item \textbf{Case study:} Walk-forward $+$ cross-sectional demean $+$ cost adjustment
        is the minimal honest backtest; leakage is the silent killer.
\end{itemize}

\medskip
\begin{block}{Next practical}
Practical 11: regulatory \& compliance applications --- structured extraction, auditability,
and SR 11-7 model-risk obligations for LLM outputs.
\end{block}
\end{frame}

% ============================================================
\appendixstart{Stretch Problems}
% ============================================================

\begin{frame}[noframenumbering]{Appendix: Stretch Problem A --- CVaR Budget}
A position has daily returns approximated by a fat-tailed distribution with
$\text{VaR}_{0.05}=-3.0\%$ and an expected shortfall (CVaR) of $-4.5\%$ in the worst 5\%.

\medskip
\begin{itemize}
  \item[\textbf{A.}] Interpret both numbers in one sentence each.
  \item[\textbf{B.}] A textual risk signal $v_{t-1}$ spikes; the GARCH update
        $\sigma_t^2=\omega+\varphi_1\varepsilon_{t-1}^2+\beta_1\sigma_{t-1}^2+\kappa v_{t-1}$
        with $\kappa>0$ raises $\sigma_t$ by 20\%. Qualitatively, what happens to
        $\text{CVaR}_{0.05}$ and to the CVaR-constrained weight on this position?
  \item[\textbf{C.}] Why is CVaR preferred to VaR for \emph{optimization}?
        (Hint: coherence; convex/LP-representable.)
\end{itemize}
\end{frame}

\begin{frame}[noframenumbering]{Appendix: Stretch Problem B --- Can You Trust the LLM Strategy's Track Record?}
Your LLM news-sentiment strategy reports a daily $\text{SR}^*=0.08$. Before allocating, ask
how long a record is needed to believe it. The minimum track record length:
\[
T^*=1+\Bigl(1-\gamma_1\text{SR}^*+\tfrac{\gamma_2-1}{4}(\text{SR}^*)^2\Bigr)\frac{Z_{1-\alpha}^2}{(\text{SR}^*)^2},
\]
with skewness $\gamma_1=-0.5$, kurtosis $\gamma_2=6$, $Z_{0.95}=1.645$:

\begin{itemize}
  \item[\textbf{A.}] Compute $T^*$ in trading days, then in years ($\approx252$/yr).
  \item[\textbf{B.}] Repeat for a symmetric, normal benchmark ($\gamma_1=0,\gamma_2=3$).
        How much does the fat left tail inflate the required record?
  \item[\textbf{C.}] The LLM strategy has only 18 months of genuine out-of-sample data
        (the model was pre-trained on text overlapping earlier periods). Given your $T^*$,
        is that enough to deploy capital? Connect to the temporal-leakage concern for
        text-trained signals.
\end{itemize}
\end{frame}

\end{document}
